% ============================================
% LUC-10
% Gerado de molambudos.md — build_miolo.py
% ============================================
\section*{\textbf{LUC-10} --- O Laudo de Lúcia}
\markboth{LUC-10 --- O Laudo de Lúcia}{}

\noindent \textit{O autorretrato diagnóstico}

\noindent \textbf{Julho de 2026 --- Belo Horizonte}
\textbf{Último registro profissional de Dra. Lúcia Mendes}

\vspace{0.5em}\noindent\rule{\textwidth}{0.5pt}\vspace{0.5em}

\noindent \textbf{Parecer Psiquiátrico --- Paciente: Lúcia Mendes}
\textbf{CRM-MG 28.391}
\textbf{Data: 28 de julho de 2026}

\vspace{0.5em}\noindent\rule{\textwidth}{0.5pt}\vspace{0.5em}

\noindent Estou escrevendo este laudo sobre mim mesma.

\noindent Sei que não é ético. Mas não há outro psiquiatra que possa avaliar meu caso. Porque meu caso não é comum. Meu caso é uma contaminação.

\noindent \textbf{Histórico:}

\noindent Em janeiro de 2026, entrei em contato com o diário do paciente 1.260 (Joaquim Antônio Correia) para fins de pesquisa acadêmica. O diário descreve o encontro de Joaquim com uma entidade não identificada (\textquotedblleft{}a criatura\textquotedblright{}) em uma vala de corpos no Campo do Alagadiço, próximo a Fortaleza, em 1915.

\noindent Após a leitura, desenvolvi:

\noindent 1. \textbf{Sintomas somáticos}: Olho direito com alteração de coloração para amarelo-dourado. Relato de \textquotedblleft{}cheiro adocicado\textquotedblright{} constante. Fadiga crônica.

\noindent 2. \textbf{Sintomas psíquicos}: Pesadelos recorrentes com a vala. Pensamentos intrusivos sobre a criatura. Sensação de presença no quarto. Hipervigilância. Desrealização.

\noindent 3. \textbf{Sintomas comportamentais}: Isolamento social. Dificuldade de concentração. Necessidade compulsiva de reler o diário.

\noindent \textbf{Diagnóstico:}

\noindent - \textbf{F43.10 --- TEPT Complexo}: Preencho todos os critérios (A a G). A exposição foi a leitura do diário. A reexperiência ocorre em sonhos e flashbacks diurnos. A evitação é inútil --- o diário me segue.

\noindent - \textbf{F44.81 --- Transtorno Dissociativo Leve}: Sinto que há uma \textquotedblleft{}segunda voz\textquotedblright{\textit{ dentro de mim. Não é alucinação auditiva --- é uma }presença\textit{ que comenta, observa, espera. A voz diz: }\textquotedblleft{}Você é minha agora.\textquotedblright{}}

\noindent - \textbf{Z65.8 --- Circunstância de Privação Extrema}: Fui privada da minha sanidade por um texto. A literatura nunca me pareceu perigosa. Agora sei que algumas palavras são armadilhas.

\noindent \textbf{Considerações Finais:}

\noindent Emito este laudo para registro. Não como profissional --- como vítima.

\noindent O caso do paciente 1.260 não é um caso psiquiátrico. É um caso de \textit{transmissão}. A criatura não é um delírio. É um parasita narrativo. Ela se transmite pela leitura.

\noindent Quem lê, recebe.
Quem recebe, carrega.
Quem carrega, passa adiante.

\noindent Eu li.
Eu recebi.
Eu carrego.

\noindent Agora preciso passar adiante.

\noindent Para você.

\vspace{0.5em}\noindent\rule{\textwidth}{0.5pt}\vspace{0.5em}

\noindent \textit{--- Dra. Lúcia Mendes}
\textit{--- Psiquiatra Forense}
\textit{--- Paciente 1.262}
\textit{--- Próxima: você}

\vspace{0.5em}
\noindent\textit{\textbf{↪ Rotas:} \rota{LUC-11} \rotasep \rota{DOC-08}}